\scanpage{87}\mspage{T-1}

\begin{center}
COURBES ELLIPTIQUES - FORMULAIRE

(Traduction libre d'un manuscrit inédit de TATE)
\end{center}

\figure{\begin{tikzpicture}[x=0.016cm,y=0.016cm]
\draw[gray,thin] (0,0) .. controls (8,0) and (13,9) .. (22,14) .. controls (35,29) and (56,42) .. (77,43) .. controls (93,43) and (95,46) .. (97,50) -- (101,46);
\draw[gray,thin] (25,15) .. controls (32,0) and (43,-13) .. (47,-27);
\end{tikzpicture}}

1. \uline{Formules générales}

Si l'on part d'une c\deleted{\unsure{a|u}}ourbe elliptique $E$, munie d'une origine $O$, et d'une forme de première espèce $\pi\ne 0$, on peut trouver $x$ et $y$ sur $E$, telles que $x$ (resp. $y$) ait un pôle double (resp. triple) en $O$, et soit morphique ailleurs; de plus, si $\pi=dt+\ldots$, on peut s'arranger (quitte à multiplier $x$ et $y$ par des constantes) pour que $x\handwritten{=}{t}^{-2}+\ldots$, $y=-t^{-3}+\ldots$

Dans le plongement projectif défini par le diviseur $3(O)$, l'équation de la courbe $E$ s'écrit :
{\makeatletter\tagsleft@true\makeatother
\[
 y^{2}+a_{1}xy+a_{3}y=x^{3}+a_{2}x^{2}+a_{4}x+a_{6}
 \tag*{$(\ast)$}
\]
}

(Homogénéite : $y$ est de poids $3$, $x$ de poids $2$, et les $a_{i}$ de poids $i$. Cela signifie que, si l'on remplace $\pi$ par $u\pi$, $x$ est remplacé par $u^{-2}x$, $y$ par $u^{-3}y$, etc.)

Inversement, partons d'une équation de la forme $(\ast)$. On posera :
\[
\begin{aligned}
& b_{2}=a_{1}^{2}+4a_{2} \qquad b_{4}=a_{1}a_{3}+2a_{4} \qquad b_{6}=a_{3}^{2}+4a_{6} \\
& -b_{8}=a_{1}a_{3}a_{4}+a_{4}^{2}-a_{1}^{2}a_{6}-a_{2}a_{3}^{2}-4a_{2}a_{6} \\
& b_{8}=b_{2}a_{6}+a_{2}a_{3}^{2}-a_{1}a_{3}a_{4}-a_{4}^{2}=a_{2}b_{6}+a_{1}^{2}a_{6}-a_{4}^{2}-a_{1}a_{3}a_{4} \\
& \hspace{4em}-4\,b_{8}=b_{4}^{2}-b_{2}b_{6}
\end{aligned}
\]
% archive sheet number: 86
